% GENERATED FILE. Edit canonical lessons, then run npm run generate:books.
% canonical-source-hash: fnv1a64:6168797b65f9cb74
% canonical-lessons: LA-C532-manumissio, LA-C532-libertinus, LA-C532-vindicta, LA-C532-ingenuus, LA-C532-pilleus

\chapter{Freedom}
\label{ch:la-freedom}
\hlchaptermodality{532}

\begin{chapteropening}
\textbf{By the end of this chapter:} \emph{I can talk about freeing slaves and freed people.}

Complete the last lesson of chapter 532: I can talk about freeing slaves and freed people.
\end{chapteropening}

\section[manūmissiō, manūmissiōnis]{manūmissiō, manūmissiōnis — the freeing of a slave}

\label{lesson:LA-C532-manumissio}



\begin{quote}
Before the new one: say the Latin for a litter-bearer, then the Latin for a farm housekeeper.
\end{quote}

\subsection*{You'll want to know: manūmissiō, manūmissiōnis}
\textbf{manūmissiō, manūmissiōnis} — \textquotedblleft{}the freeing of a slave\textquotedblright{}.

\begin{quote}
Post manūmissiōnem servus lībertus est.
\end{quote}

\begin{quote}
\emph{After being set free, the slave is a freedman.}
\end{quote}

Feminine, third declension: \emph{manus} 'hand' + \emph{mittō} 'to send' — to let go from the hand. English \emph{manumission}.

\begin{cousinweb}
\textbf{manūmissiō} is built on \emph{manūmittō} \textquotedblleft{}to set free\textquotedblright{}, from \emph{manus} \textquotedblleft{}hand, power\textquotedblright{} and \emph{mittō} \textquotedblleft{}to let go\textquotedblright{} (Lewis \& Short). English \emph{manumission} comes from it.
\end{cousinweb}

\begin{grammarlens}[title={What you've built}]
One more word for telling, explaining and describing at B1.
\end{grammarlens}

\subsection*{Your turn}
\begin{itemize}
\raggedright
  \item \emph{Say it:} \emph{manūmissiō, manūmissiōnis}
  \item \emph{Say it:} \emph{manūmissiō, manūmissiōnis}, once more
  \item \emph{Say it:} a sentence of your own with \emph{manūmissiō, manūmissiōnis}
  \item \emph{Recall:} say the Latin for a litter-bearer, then the Latin for a farm housekeeper, then say \emph{manūmissiō, manūmissiōnis} again
\end{itemize}

\begin{tcolorbox}[breakable,colback=teal!4,colframe=teal!35!black,title={Before you move on}]
What does manūmissiō, manūmissiōnis mean? (\textquotedblleft{}The freeing of a slave\textquotedblright{}.) Say it once more.
\end{tcolorbox}
\section[lībertīnus, lībertīnī]{lībertīnus, lībertīnī — a freedman (by rank)}

\label{lesson:LA-C532-libertinus}



\begin{quote}
Before the new one: say the Latin for a farm housekeeper, then the Latin for the freeing of a slave.
\end{quote}

\subsection*{You'll want to know: lībertīnus, lībertīnī}
\textbf{lībertīnus, lībertīnī} — \textquotedblleft{}a freedman (by rank)\textquotedblright{}.

\begin{quote}
Multī lībertīnī in urbe tabernās habent.
\end{quote}

\begin{quote}
\emph{Many freedmen keep shops in the city.}
\end{quote}

Masculine. \emph{Lībertīnus} names a freed slave's rank in society, \emph{lībertus} his tie to his former master.

\begin{cousinweb}
\textbf{lībertīnus} is built on \emph{lībertus} \textquotedblleft{}freedman\textquotedblright{} (Lewis \& Short; EtymDB), from \emph{līber} \textquotedblleft{}free\textquotedblright{}.
\end{cousinweb}

\begin{grammarlens}[title={What you've built}]
One more word for telling, explaining and describing at B1.
\end{grammarlens}

\subsection*{Your turn}
\begin{itemize}
\raggedright
  \item \emph{Say it:} \emph{lībertīnus, lībertīnī}
  \item \emph{Say it:} \emph{lībertīnus, lībertīnī}, once more
  \item \emph{Say it:} a sentence of your own with \emph{lībertīnus, lībertīnī}
  \item \emph{Recall:} say the Latin for a farm housekeeper, then the Latin for the freeing of a slave, then say \emph{lībertīnus, lībertīnī} again
\end{itemize}

\begin{tcolorbox}[breakable,colback=teal!4,colframe=teal!35!black,title={Before you move on}]
What does lībertīnus, lībertīnī mean? (\textquotedblleft{}A freedman (by rank)\textquotedblright{}.) Say it once more.
\end{tcolorbox}
\section[vindicta, vindictae]{vindicta, vindictae — the rod of freedom (in freeing a slave); vengeance}

\label{lesson:LA-C532-vindicta}



\begin{quote}
Before the new one: say the Latin for the freeing of a slave, then the Latin for a freedman.
\end{quote}

\subsection*{You'll want to know: vindicta, vindictae}
\textbf{vindicta, vindictae} — \textquotedblleft{}the rod of freedom (in freeing a slave); vengeance\textquotedblright{}.

\begin{quote}
Magistrātus servum vindictā tangit, et servus līber fit.
\end{quote}

\begin{quote}
\emph{The magistrate touches the slave with the rod, and the slave becomes free.}
\end{quote}

Feminine, first declension. In one way of freeing a slave, he was touched with this rod before a magistrate. Later it also means vengeance; English \emph{vindictive}.

\begin{cousinweb}
\textbf{vindicta} is built on \emph{vindicō} \textquotedblleft{}to lay claim to, set free\textquotedblright{} (Lewis \& Short; EtymDB), from \emph{vindex} \textquotedblleft{}a defender\textquotedblright{}. English \emph{vindictive} comes from its later sense, vengeance.
\end{cousinweb}

\begin{grammarlens}[title={What you've built}]
One more word for telling, explaining and describing at B1.
\end{grammarlens}

\subsection*{Your turn}
\begin{itemize}
\raggedright
  \item \emph{Say it:} \emph{vindicta, vindictae}
  \item \emph{Say it:} \emph{vindicta, vindictae}, once more
  \item \emph{Say it:} a sentence of your own with \emph{vindicta, vindictae}
  \item \emph{Recall:} say the Latin for the freeing of a slave, then the Latin for a freedman, then say \emph{vindicta, vindictae} again
\end{itemize}

\begin{tcolorbox}[breakable,colback=teal!4,colframe=teal!35!black,title={Before you move on}]
What does vindicta, vindictae mean? (\textquotedblleft{}The rod of freedom (in freeing a slave); vengeance\textquotedblright{}.) Say it once more.
\end{tcolorbox}
\section[ingenuus, ingenua, ingenuum]{ingenuus, ingenua, ingenuum — free-born; frank}

\label{lesson:LA-C532-ingenuus}



\begin{quote}
Before the new one: say the Latin for a freedman, then the Latin for the rod of freedom.
\end{quote}

\subsection*{You'll want to know: ingenuus, ingenua, ingenuum}
\textbf{ingenuus, ingenua, ingenuum} — \textquotedblleft{}free-born; frank\textquotedblright{}.

\begin{quote}
Hic puer ingenuus est, nōn servus.
\end{quote}

\begin{quote}
\emph{This boy is free-born, not a slave.}
\end{quote}

Born free, as against \emph{lībertīnus}; then honest, frank. English \emph{ingenuous}.

\begin{cousinweb}
\textbf{ingenuus} is built on \emph{in-} and the root of \emph{gignō} \textquotedblleft{}to beget\textquotedblright{} (Lewis \& Short): born in the place, then born free. English \emph{ingenuous} comes from it.
\end{cousinweb}

\begin{grammarlens}[title={What you've built}]
One more word for telling, explaining and describing at B1.
\end{grammarlens}

\subsection*{Your turn}
\begin{itemize}
\raggedright
  \item \emph{Say it:} \emph{ingenuus, ingenua, ingenuum}
  \item \emph{Say it:} \emph{ingenuus, ingenua, ingenuum}, once more
  \item \emph{Say it:} a sentence of your own with \emph{ingenuus, ingenua, ingenuum}
  \item \emph{Recall:} say the Latin for a freedman, then the Latin for the rod of freedom, then say \emph{ingenuus, ingenua, ingenuum} again
\end{itemize}

\begin{tcolorbox}[breakable,colback=teal!4,colframe=teal!35!black,title={Before you move on}]
What does ingenuus, ingenua, ingenuum mean? (\textquotedblleft{}Free-born; frank\textquotedblright{}.) Say it once more.
\end{tcolorbox}
\section[pilleus, pilleī]{pilleus, pilleī — a felt cap}

\label{lesson:LA-C532-pilleus}



\begin{quote}
Before the new one: say the Latin for the rod of freedom, then the Latin for free-born.
\end{quote}

\subsection*{You'll want to know: pilleus, pilleī}
\textbf{pilleus, pilleī} — \textquotedblleft{}a felt cap\textquotedblright{}.

\begin{quote}
Servus līberātus pilleum gerit.
\end{quote}

\begin{quote}
\emph{The freed slave wears the felt cap.}
\end{quote}

Masculine, second declension. A newly freed slave put on the \emph{pilleus}, a sign of liberty.

\begin{cousinweb}
\textbf{pilleus} is usually linked with Greek \emph{pilos} \textquotedblleft{}felt, felt cap\textquotedblright{} (Lewis \& Short), but whether it is a loan or a cousin is not settled.
\end{cousinweb}

\begin{grammarlens}[title={What you've built}]
One more word for telling, explaining and describing at B1.
\end{grammarlens}

\subsection*{Your turn}
\begin{itemize}
\raggedright
  \item \emph{Say it:} \emph{pilleus, pilleī}
  \item \emph{Say it:} \emph{pilleus, pilleī}, once more
  \item \emph{Say it:} a sentence of your own with \emph{pilleus, pilleī}
  \item \emph{Recall:} say the Latin for the rod of freedom, then the Latin for free-born, then say \emph{pilleus, pilleī} again
\end{itemize}

\begin{tcolorbox}[breakable,colback=teal!4,colframe=teal!35!black,title={Before you move on}]
What does pilleus, pilleī mean? (\textquotedblleft{}A felt cap\textquotedblright{}.) Say it once more.
\end{tcolorbox}
